% Lösungen Einheit 1 von 3 – Bestand aus Rate (zusammenbau v0.1)
\einheitenkopf{Einheit 1 von 3 · Bestand aus Rate}

\erg{8}{a) eine Rate – Liter je Minute ist eine Menge pro Zeit \quad b) $[t^3 - 6t^2 + 15t]_0^4 = 28$; in den ersten vier Stunden fließen 28 Liter zu \quad c) $50 + ∫_0^3 r(t)\,\mathrm{d}t = 50 + [2t^2 + 2t]_0^3 = 50 + 24 = 74$ Liter \quad d) Änderung $\approx 4$\,°C (Fläche über der Achse von 0 bis 4); die Rate ist dort positiv, die Temperatur nimmt zu \quad e) $∫_0^2 90t\,\mathrm{d}t = 180$ (km/h mal Minuten); geteilt durch 60: Strecke $= 3\,\mathrm{km}$ \quad f) $5 + ∫_0^T r(t)\,\mathrm{d}t = 20$ \quad g) $∫_1^5 r(t)\,\mathrm{d}t = [300t + 30t^2 - 2t^3]_1^5 = 2\,000 - 328 = 1\,672$; Liste: $2\,600 - 1\,000 = 1\,600$; Abweichung $= 72 : 1\,600 = 4{,}5\,\%$}
\erg{9}{a) $∫_0^T r(t)\,\mathrm{d}t = T^3 - 6T^2 = T^2 \cdot (T - 6) = 0$, also $T = 6$}
\erg{10}{a) Jonas vergisst den Anfangsbestand; das Integral ist nur die Zunahme. Inhalt $= 80 + 48 = 128$ Liter \quad b) Das Integral der Rate summiert alle kleinen Zuwächse Rate mal Zeit, also die gesamte Änderung; zum Anfangsbestand addiert ergibt das den Endbestand. \quad c) Fläche zwischen dem Graphen von r und der t-Achse von t = 1 bis t = 3; sie gibt an, wie viele Liter zwischen Minute 1 und Minute 3 zufließen \quad d) $∫_0^6 r(t)\,\mathrm{d}t = 1\,440 - 360 = 1\,080\,\mathrm{m}^3 > 900\,\mathrm{m}^3$ – ja, es läuft über}
